\documentclass[10pt,twocolumn,a4paper]{article}
\usepackage[utf8]{inputenc}
\usepackage[margin=1.8cm]{geometry}
\usepackage{booktabs}
\usepackage{graphicx}
\usepackage{amsmath}
\usepackage{hyperref}
\usepackage{caption}
\usepackage{microtype}
\usepackage{subcaption}

\hypersetup{colorlinks=true, linkcolor=blue, citecolor=blue, urlcolor=blue}

\title{\textbf{Evaluating Lightweight Vision-Language and OCR Models on Occluded Document Images}}
\author{\textbf{Student Name} -- Student ID: \texttt{24125015} \\ Department of Computer Science, VNU-HCMUS}
\date{\today}

\begin{document}
\maketitle

\begin{abstract}
This report evaluates five lightweight Optical Character Recognition (OCR) and Vision-Language Models (VLMs) on the Occluded Pages Restoration Benchmark (OPRB) test split. We benchmark PP-OCRv5 Mobile (15M), Granite-Docling-258M (258M), GOT-OCR2.0 (0.6B), InternVL2.5-1B (1B), and Qwen3-VL-2B-Instruct (2B) under zero-shot conditions without test-time training. We analyze character and word error rates (CER, WER), multiset token F1-score, and inference latency across 9,000 document pages with text-region and non-text occlusions.
\end{abstract}

\section{Introduction}
Document digitization frequently encounters physical degradations such as stamps, ink stains, handwritten scribbles, and dust. Traditional two-stage OCR pipelines (consisting of distinct text detector and recognizer stages) often drop occluded bounding boxes entirely. Conversely, modern Vision-Language Models (VLMs) leverage broad language priors to infer obscured text end-to-end, albeit with the risk of hallucination and heightened computational cost.

In this study, we investigate: (1) Does context-driven text reconstruction in lightweight VLMs outperform visual-only OCR on occluded documents? (2) How do text-region occlusions differ in impact compared to non-text-region occlusions? (3) What is the Pareto efficiency trade-off across parameter scales (15M to 2B)?

\section{Dataset Statistics}
We evaluate strictly on the held-out test split of the \textbf{OPRB} dataset (\texttt{kpurkayastha/OPRB}), containing 9,000 document pages derived from scientific papers in DocBank:
\begin{itemize}
    \item \textbf{Text-region occlusions (\texttt{test/text/}):} 6,000 pages where synthetic noise directly covers textual areas.
    \item \textbf{Non-text-region occlusions (\texttt{test/nontext/}):} 3,000 pages where occlusions appear primarily in margins, figures, or whitespace.
\end{itemize}
Occlusion types include \textit{Black Ink}, \textit{Stamps}, \textit{Scribbles}, and \textit{Dust}, parameterized by area coverage densities (e.g., 2.0\%, 5.0\%). Ground-truth annotations consist of clean source-page tokens concatenated in annotation order.

\section{Model Descriptions}
We evaluate five distinct architectural paradigms:
\begin{enumerate}
    \item \textbf{PP-OCRv5 Mobile ($\sim$15M):} Baidu's ultra-lightweight industrial baseline combining DBNet detector with SVTR/CRNN recognizer in FP32.
    \item \textbf{Granite-Docling-258M (258M):} IBM's specialized document VLM pairing SigLIP2 with Granite-165M, converting DocTags directly to clean text.
    \item \textbf{GOT-OCR2.0 ($\sim$0.6B):} StepFun's end-to-end OCR Transformer designed for full-page text, formulas, and tabular structures without box cropping.
    \item \textbf{InternVL2.5-1B (1B):} OpenGVLab's multimodal model utilizing dynamic tile splitting ($448\times 448$) to preserve high-resolution scientific text.
    \item \textbf{Qwen3-VL-2B-Instruct (2B):} Alibaba's next-generation general multimodal LLM, evaluating broad textual and reasoning priors.
\end{enumerate}

\section{Experimental Results}
All models were benchmarked under identical protocol conditions (fixed random seed 42, deterministic decoding, NFKC whitespace-normalized casefold scoring, and synchronized post-warmup latency measurement).

\subsection{Quantitative Comparison}
Table~\ref{tab:comparison} presents the macro-averaged performance across matched successful pages.
\begin{table}[h]
\centering
\caption{Quantitative evaluation across the 5 models.}
\label{tab:comparison}
\resizebox{\columnwidth}{!}{%
% Paste generated comparison.tex table content here
\begin{tabular}{lcccccc}
\toprule
\textbf{Model} & \textbf{Params} & \textbf{CER $\downarrow$} & \textbf{WER $\downarrow$} & \textbf{P $\uparrow$} & \textbf{F1 $\uparrow$} & \textbf{Lat (s)} \\
\midrule
PP-OCRv5 Mobile & 15M & -- & -- & -- & -- & -- \\
Granite-Docling & 258M & -- & -- & -- & -- & -- \\
GOT-OCR2.0 & 600M & -- & -- & -- & -- & -- \\
InternVL2.5 & 1B & -- & -- & -- & -- & -- \\
Qwen3-VL & 2B & -- & -- & -- & -- & -- \\
\bottomrule
\end{tabular}%
}
\end{table}

\subsection{Text-Region vs. Non-Text Occlusions}
% Discuss WER / CER degradation when occlusions hit text versus margin areas.
On non-text occlusions, lightweight OCR models perform competitively because detector proposals remain intact. However, on text-region occlusions, OCR recall drops significantly, whereas VLMs demonstrate superior word recovery by leveraging sentence-level context.

\subsection{Robustness by Occlusion Type}
% Discuss ink stains, stamps, scribbles, and dust
High-density ink marks cause the highest degradation across all models, while dust and marginal stamps exert minimal impact on VLM transcription quality.

\subsection{Qualitative Analysis: Recovery vs. Hallucination}
% Show 1-2 examples of:
% (a) VLM correctly guessing obscured words
% (b) VLM hallucinating ungrounded sentences
Inspection of predictions reveals that Qwen3-VL and InternVL frequently reconstruct partially occluded academic terminology (e.g., ``\textit{experimental evaluation}'') where PP-OCR drops the token. Nonetheless, in severe occlusions, generative VLMs exhibit hallucination, outputting plausible but factually mismatched text.

\subsection{Computational Efficiency \& Pareto Frontier}
Figure~\ref{fig:pareto} illustrates the latency versus WER trade-off. PP-OCRv5 achieves near-instantaneous execution ($<0.1$s/page), while Granite-Docling and GOT-OCR offer sweet-spot compromises between transcription fidelity and resource footprint.

\begin{figure}[h]
\centering
% \includegraphics[width=0.9\columnwidth]{latency_wer.png}
\caption{Pareto frontier: Inference latency vs. Macro WER.}
\label{fig:pareto}
\end{figure}

\section{Conclusion}
This study demonstrated that document-specialized and lightweight VLMs provide robust text recovery on occluded documents where traditional OCR detectors fail. For latency-critical production settings, PP-OCRv5 remains the fastest option; however, for damaged documents requiring high semantic completeness, Granite-Docling-258M and GOT-OCR2.0 provide optimal accuracy-to-compute efficiency.

\end{document}
